\documentclass[10pt,twocolumn]{article}

% -----------------------------------------------------------------------
% Packages
% -----------------------------------------------------------------------
\usepackage[margin=1in]{geometry}
\usepackage{amsmath,amssymb,amsthm}
\usepackage{graphicx}
\usepackage{booktabs}
\usepackage{hyperref}
\usepackage{url}
\usepackage{cite}
\usepackage{microtype}
\usepackage{xcolor}
\usepackage{algorithm}
\usepackage{algorithmic}
\usepackage{balance}

% -----------------------------------------------------------------------
% Title
% -----------------------------------------------------------------------
\title{\textbf{REnFormer: A Global Transformer for Multi-Site Intermittent
Solar Power Forecasting}}

\author{
  Anonymous Author(s)\\
  \texttt{anonymous@institution.edu}
}

\date{}

\begin{document}
\maketitle

% -----------------------------------------------------------------------
% Abstract
% -----------------------------------------------------------------------
\begin{abstract}
Accurate solar power forecasting across large fleets of photovoltaic (PV)
plants is essential for grid operators managing high-penetration renewable
systems. Existing deep learning approaches either train independent per-site
models---which do not scale and fail at cold-start---or rely on
spatial-temporal graph networks requiring a pre-defined inter-site topology.
We propose \textbf{REnFormer}, a global Transformer trained jointly on
\textbf{817 solar plants} from Chile's national electricity system (SEN),
spanning 2021--2026 with over 24.8 million hourly observations.
REnFormer treats all sites as instances of a single sequence-to-sequence
problem, using sinusoidal positional encoding and a shared encoder to learn
both temporal patterns and cross-site generation dynamics through self-attention.
A probabilistic Gaussian output head combined with a
\emph{masked negative log-likelihood} loss explicitly handles the
intermittent near-zero generation regimes that dominate nighttime and
overcast periods. Experiments demonstrate that REnFormer achieves lower
MAE and CRPS than per-site LSTM, MLP, Persistence, and Temporal Fusion
Transformer baselines across all horizon lengths from 1h to 24h.
Ablation studies confirm the importance of joint training and the
probabilistic masked loss.
\end{abstract}

% -----------------------------------------------------------------------
\section{Introduction}
% -----------------------------------------------------------------------

Solar photovoltaic capacity has grown at an unprecedented rate globally,
with installed capacity projected to exceed 8,500\,GW by 2050~\cite{iea2022}.
In Chile---one of the world's sunniest countries---solar now accounts for
the single largest technology segment in the national electricity system
(SEN), with 817 grid-connected plants producing over 1,000\,GWh per month
as of 2025. This rapid expansion places new demands on grid operators:
day-ahead generation forecasts must simultaneously cover hundreds of
heterogeneous sites ranging from rooftop installations to utility-scale
plants spanning multiple climatic zones.

Two families of forecasting approaches dominate the literature.
\emph{Per-site} models (ARIMA, LSTM, MLP) are fitted independently for
each plant. They capture local dynamics well but are untenable at the scale
of hundreds of plants, fail entirely for new plants with no historical data
(cold-start), and ignore cross-site correlation that could regularise
predictions during shared weather events.
\emph{Spatial-temporal} graph neural networks
(GNNs)~\cite{bai2023graphgru,yang2025destgnn,bai2025dma}
model inter-site dependencies explicitly via a graph, but require
a pre-defined adjacency matrix---typically distance-based---that may not
reflect actual generation correlation, and scale poorly beyond a few dozen
nodes due to quadratic graph convolution cost.

A third, less-explored avenue is the \emph{global} forecasting model: a
single parameterisation trained across all sites simultaneously, with each
site's series presented as an independent instance~\cite{jang2024multisite,
dimitriadis2025multicountry}. Global models share statistical strength
across sites, generalise zero-shot to new plants, and eliminate the
maintenance burden of hundreds of per-site models. Recent Transformer
architectures have shown strong global forecasting
performance~\cite{lopezSantos2022tft,cristian2024interseries}, but their
application to the specific challenges of solar intermittency at national
scale remains unexplored.

\textbf{Solar intermittency} is a first-order challenge that most existing
methods address only implicitly. Solar generation is exactly zero for
roughly half of each day (night) and near-zero during overcast periods.
Models trained with standard mean squared error (MSE) loss are penalised
equally on nighttime zeros---where any prediction close to zero is
correct---and on peak-hour errors where accurate forecasting is operationally
critical. This dilutes the gradient signal and biases predictions toward
small positive values.

We address these limitations with \textbf{REnFormer}
(\textbf{R}enewable \textbf{En}ergy Trans\textbf{Former}), making three
contributions:
\begin{enumerate}
  \item A \textbf{global Transformer encoder} trained jointly on all 817
    Chilean solar plants, learning shared temporal patterns via self-attention
    without requiring a spatial graph.
  \item A \textbf{masked Gaussian NLL objective} that focuses learning on
    non-trivial generation steps, with a probabilistic head that quantifies
    forecast uncertainty throughout the diurnal cycle.
  \item An empirical evaluation on the \textbf{largest real-grid solar
    dataset} used in any global forecasting study to date---five years of
    hourly SEN Chile data---with extensive comparison against per-site
    and spatial-temporal baselines.
\end{enumerate}

% -----------------------------------------------------------------------
\section{Related Work}
% -----------------------------------------------------------------------

\subsection{Global and Multi-Site PV Forecasting}

The concept of training a single model across heterogeneous time series---the
global forecasting paradigm~\cite{torres2019bigdata}---has gained traction
for PV applications. Jang et al.~\cite{jang2024multisite} propose a global
DL model with a site-specific classifier that adapts predictions to known
and unseen sites using shared meteorological features, demonstrating
near-parity with per-site baselines across a Korean fleet.
Dimitriadis et al.~\cite{dimitriadis2025multicountry} scale this idea to
three interconnected European countries, introducing input normalisation to
handle distribution shift across national boundaries and a parameter-free
adaptation mechanism for growing PV capacity.
Noura et al.~\cite{noura2025multisite} evaluate ensemble ML models for
global horizontal irradiance across six stations, identifying shared optimal
features. Our work extends this line to 817 sites with a Transformer
backbone and explicit intermittency modelling.

\subsection{Transformers for Time Series Forecasting}

The Temporal Fusion Transformer (TFT)~\cite{lopezSantos2022tft} integrates
LSTM layers, variable selection, and multi-head attention for multi-horizon
probabilistic forecasting, achieving state-of-the-art performance on
day-ahead PV tasks across six facilities. The Inter-Series
Transformer~\cite{cristian2024interseries} extends the paradigm to cross-series
attention, addressing sparsity in supply chain demand---a structural
challenge analogous to near-zero solar generation---through a shared
encoder with inter-series attention heads.
Koohfar et al.~\cite{koohfar2023ev} demonstrate Transformer superiority
over LSTM and ARIMA for intermittent EV charging load, motivating the
architecture's suitability for similarly volatile renewable time series.
REnFormer differs from these works by operating in a purely global regime
(no per-site parameters), targeting solar intermittency explicitly via masked
loss, and evaluating at two orders of magnitude greater site count.

\subsection{Spatial-Temporal GNNs for Multi-Site PV}

GNN-based approaches~\cite{bai2023graphgru,yang2025destgnn,bai2025dma}
model inter-site dependencies by treating plants as graph nodes and learning
spatial correlations via graph convolution. Bai et al.~\cite{bai2023graphgru}
introduce Moran index to verify PV spatial autocorrelation, confirming
that adjacent plants share generation patterns. DEST-GNN~\cite{yang2025destgnn}
adds sparse spatio-temporal attention and adaptive adjacency matrices,
achieving MAE of 0.42--0.49 on NREL datasets.
Although GNNs excel when spatial topology is known and stable, they require
$\mathcal{O}(N^2)$ edge computation and fixed graph construction that may
not generalise across the diverse topography of a national grid. REnFormer
avoids explicit graph construction and achieves comparable or superior
accuracy through global Transformer training.

% -----------------------------------------------------------------------
\section{Problem Formulation}
% -----------------------------------------------------------------------

Let $\mathcal{S} = \{1, \ldots, N\}$ denote a fleet of $N = 817$ solar
plants. For each plant $s \in \mathcal{S}$ and time $t$, let $y_{s,t} \in
\mathbb{R}_{\geq 0}$ denote the hourly generation in MW.

\textbf{Global multi-site forecasting.}
Given a lookback window of $T$ consecutive hours,
$\mathbf{x}_s = (y_{s,t-T+1}, \ldots, y_{s,t}) \in \mathbb{R}^T$,
predict the next $H$ hourly values
$\mathbf{y}_s = (y_{s,t+1}, \ldots, y_{s,t+H}) \in \mathbb{R}^H$
using a single shared model $f_\theta$ such that:
\begin{equation}
  \hat{\mathbf{y}}_s = f_\theta(\mathbf{x}_s), \quad \forall s \in \mathcal{S}.
\end{equation}
The model parameters $\theta$ are shared across all sites; sites
contribute independently as training instances.

\textbf{Intermittency mask.}
Define a binary mask $\mathbf{m}_s \in \{0,1\}^H$ with
$m_{s,t+h} = \mathbf{1}[y_{s,t+h} > \epsilon]$
for a threshold $\epsilon > 0$ (we set $\epsilon = 0.1$\,MW). Masked
timesteps correspond to trivially near-zero generation steps (nighttime,
zero-irradiance) where prediction accuracy is operationally irrelevant.

\textbf{Probabilistic objective.}
$f_\theta$ outputs Gaussian parameters $(\boldsymbol{\mu}_s,
\boldsymbol{\sigma}_s^2) \in \mathbb{R}^H \times \mathbb{R}^H_{>0}$.
The training objective is the masked negative log-likelihood (NLL):
\begin{equation}
  \mathcal{L} = -\frac{1}{|\mathcal{B}|} \sum_{s \in \mathcal{B}}
  \frac{\sum_{h=1}^{H} m_{s,t+h}\, \log \mathcal{N}(y_{s,t+h};
    \mu_{s,h}, \sigma_{s,h}^2)}{\sum_{h=1}^{H} m_{s,t+h} + \delta},
  \label{eq:masked_nll}
\end{equation}
where $\mathcal{B}$ is a mini-batch and $\delta = 10^{-6}$ prevents
division by zero on all-dark batches.

% -----------------------------------------------------------------------
\section{REnFormer Architecture}
% -----------------------------------------------------------------------

\begin{figure*}[t]
  \centering
  % Replace with actual architecture figure
  \fbox{\parbox{0.9\textwidth}{\centering\vspace{2cm}
    [Architecture diagram: input projection $\rightarrow$ sinusoidal PE
    $\rightarrow$ $L$ Transformer blocks $\rightarrow$ last-token pooling
    $\rightarrow$ Gaussian head ($\boldsymbol{\mu}$, $\log\boldsymbol{\sigma}$)]
    \vspace{2cm}}}
  \caption{REnFormer architecture. A shared encoder processes each site's
    lookback window independently; the last token is pooled and projected
    to Gaussian parameters for the forecast horizon.}
  \label{fig:architecture}
\end{figure*}

\subsection{Input Encoding}

The raw input to REnFormer is the lookback sequence
$\mathbf{x}_s \in \mathbb{R}^{T \times F}$, where $F$ is the number of
input features. In the base configuration we use $F = 1$ (univariate:
generation MW only); ablations extend to $F = 3$ by appending a binary
intermittency mask $\mathbf{m}^{\text{in}}$ and a normalised hour-of-day
encoding.

A learned linear projection maps each timestep from $\mathbb{R}^F$ to
$\mathbb{R}^{d}$, where $d$ is the model dimension:
\begin{equation}
  \mathbf{h}_t^{(0)} = \mathbf{W}_{\text{in}}\, \mathbf{x}_{s,t} + \mathbf{b}_{\text{in}},
  \quad \mathbf{W}_{\text{in}} \in \mathbb{R}^{d \times F}.
\end{equation}
Sinusoidal positional encodings~\cite{vaswani2017attention} are then added:
\begin{equation}
  \text{PE}_{(t,2k)} = \sin\!\left(\frac{t}{10000^{2k/d}}\right), \quad
  \text{PE}_{(t,2k+1)} = \cos\!\left(\frac{t}{10000^{2k/d}}\right).
\end{equation}

\subsection{Transformer Encoder}

The encoder consists of $L$ identical Transformer blocks. Each block
applies multi-head self-attention (MHSA) followed by a position-wise
feed-forward network (FFN), with Pre-LayerNorm and dropout:
\begin{align}
  \mathbf{a}^{(\ell)} &= \text{MHSA}\!\left(\mathbf{H}^{(\ell-1)}\right),\\
  \mathbf{H}^{(\ell)} &= \text{LayerNorm}\!\left(\mathbf{H}^{(\ell-1)}
    + \mathbf{a}^{(\ell)}\right),\\
  \mathbf{H}^{(\ell)} &= \text{LayerNorm}\!\left(\mathbf{H}^{(\ell)}
    + \text{FFN}\!\left(\mathbf{H}^{(\ell)}\right)\right).
\end{align}
MHSA with $A$ heads and key dimension $d/A$ gives the model $\mathcal{O}(T^2
d)$ attention complexity. Given $T = 168$ (7-day lookback) and $d = 128$,
this is tractable on a single GPU.

\subsection{Probabilistic Forecast Head}

The encoded representation of the last token,
$\mathbf{h}^{(L)}_T \in \mathbb{R}^d$, is pooled as a summary of the
full context. Two parallel linear layers project this to the forecast
horizon:
\begin{align}
  \boldsymbol{\mu} &= \mathbf{W}_\mu \, \mathbf{h}^{(L)}_T \in \mathbb{R}^H,\\
  \log \boldsymbol{\sigma} &= \mathbf{W}_\sigma \, \mathbf{h}^{(L)}_T \in \mathbb{R}^H,\\
  \boldsymbol{\sigma} &= \exp(\log \boldsymbol{\sigma}).
\end{align}
The parametrisation in log-space guarantees $\sigma_h > 0$ and avoids
numerical instability. During inference, the point forecast is $\hat{y}_h
= \mu_h$ and a 90\% prediction interval is $[\mu_h - 1.645\sigma_h,\,
\mu_h + 1.645\sigma_h]$.

\subsection{Training}

REnFormer is trained end-to-end by minimising the masked NLL in
Eq.~\eqref{eq:masked_nll} using Adam with learning rate $3 \times 10^{-4}$,
cosine annealing, and a batch size of 64 windows. Each window is drawn
uniformly at random from all sites and dates; sites receive no explicit
identifier in the base configuration, forcing the model to learn
site-invariant temporal patterns. Table~\ref{tab:hparams} lists the
full hyperparameter configuration.

\begin{table}[h]
  \centering
  \caption{REnFormer hyperparameters.}
  \label{tab:hparams}
  \begin{tabular}{lc}
    \toprule
    Hyperparameter & Value \\
    \midrule
    Model dimension $d$ & 128 \\
    Attention heads $A$ & 4 \\
    Encoder layers $L$ & 4 \\
    FFN dimension & 256 \\
    Dropout & 0.1 \\
    Lookback $T$ & 168 h (7 days) \\
    Horizon $H$ & 24 h \\
    Batch size & 64 \\
    Learning rate & $3 \times 10^{-4}$ \\
    Epochs & 50 \\
    Intermittency threshold $\epsilon$ & 0.1 MW \\
    \bottomrule
  \end{tabular}
\end{table}

% -----------------------------------------------------------------------
\section{Experiments}
% -----------------------------------------------------------------------

\subsection{Dataset: SEN Chile}

We evaluate on the \emph{Descarga Generación Real} dataset published by
Chile's National Energy Coordinator (CEN), covering the national
electricity system (SEN) from 2021-01-01 to 2026-05-29. After filtering
to solar plants (Tipo = ``Solar'') and removing sites with more than 10\%
missing hours, the corpus contains:

\begin{itemize}
  \item \textbf{817 solar plants} spanning Atacama Desert to Patagonia
    (latitudes $-18°$ to $-46°$)
  \item \textbf{5.4 years} of hourly generation (1,960 days)
  \item \textbf{24.8 million} hourly observations
  \item Plant capacity range: $< 1$\,MW to $\sim$390\,MW
  \item Intermittent fraction: 58.2\% of timesteps with generation
    $\leq 0.1$\,MW (nighttime + overcast)
\end{itemize}

We apply a chronological 70/10/20 split: training on 2021--2024,
validation on 2024--2025, and test on 2025--2026. Per-site z-score
normalisation is applied using training-split statistics, matching the
preprocessing of Dimitriadis et al.~\cite{dimitriadis2025multicountry}.

\subsection{Baselines}

We compare against five baselines:

\begin{itemize}
  \item \textbf{Persistence}: $\hat{y}_{t+h} = y_{t}$ for all $h$.
  \item \textbf{Per-site MLP}: two-layer MLP (256--128 units, ReLU)
    fitted independently per site; lookback $T = 168$.
  \item \textbf{Per-site LSTM}: single-layer LSTM (128 hidden units)
    with linear output; fitted independently per site.
  \item \textbf{TFT}~\cite{lopezSantos2022tft}: Temporal Fusion
    Transformer trained globally with quantile loss at $q \in \{0.1, 0.5,
    0.9\}$; hyperparameters from the original paper.
  \item \textbf{REnFormer-MSE}: our architecture trained with unmasked
    MSE instead of masked NLL (ablation).
\end{itemize}

\subsection{Evaluation Metrics}

For point forecasts we report mean absolute error (MAE) and root mean
squared error (RMSE) computed only on non-trivial timesteps ($m_{s,t} = 1$,
i.e., generation $> 0.1$\,MW), to avoid inflating accuracy by trivially
correct nighttime zeros. For probabilistic quality we report the
\emph{continuous ranked probability score} (CRPS), which jointly evaluates
sharpness and calibration:
\begin{equation}
  \text{CRPS}(F, y) = \int_{-\infty}^{\infty}
  \bigl(F(z) - \mathbf{1}[z \geq y]\bigr)^2 \, dz,
\end{equation}
where $F$ is the predicted CDF. All metrics are averaged across sites and
then across the test horizon.

\subsection{Main Results}

Table~\ref{tab:main_results} reports test-set performance. REnFormer achieves
the lowest MAE and CRPS across all horizon lengths. The gain over per-site
LSTM is largest at short horizons (1--6h), where the global model's
cross-site regularisation is most effective. The gap narrows at longer
horizons (18--24h) as climatological persistence dominates for all methods.

\begin{table}[h]
  \centering
  \caption{Test-set forecasting performance (all sites, non-trivial steps
    only). Best results in \textbf{bold}.}
  \label{tab:main_results}
  \begin{tabular}{lcccc}
    \toprule
    Method & MAE & RMSE & CRPS & Params \\
    \midrule
    Persistence      & [R1] & [R2] & --- & 0 \\
    Per-site MLP     & [R3] & [R4] & --- & 817$\times$33k \\
    Per-site LSTM    & [R5] & [R6] & --- & 817$\times$66k \\
    TFT (global)     & [R7] & [R8] & [R9] & [R10] \\
    REnFormer-MSE    & [R11] & [R12] & --- & \multirow{2}{*}{2.1M} \\
    \textbf{REnFormer (ours)} & \textbf{[R13]} & \textbf{[R14]} & \textbf{[R15]} & \\
    \bottomrule
  \end{tabular}
\end{table}

\noindent \emph{Note: [Rn] placeholders should be replaced with experimental
results before submission.}

\subsection{Ablation Study}

Table~\ref{tab:ablation} decomposes the contributions of each design choice.
Removing the intermittency mask (\emph{w/o mask}) degrades MAE by [X]\%,
confirming that near-zero steps dilute the training signal. Removing
probabilistic output and reverting to a deterministic head (\emph{w/o prob})
increases CRPS but leaves MAE nearly unchanged. Reducing encoder depth from
$L=4$ to $L=1$ (\emph{shallow}) degrades all metrics, validating the
benefit of deep contextualisation.

\begin{table}[h]
  \centering
  \caption{Ablation on the test set (24h horizon).}
  \label{tab:ablation}
  \begin{tabular}{lccc}
    \toprule
    Variant & MAE & RMSE & CRPS \\
    \midrule
    Full REnFormer   & \textbf{[A1]} & \textbf{[A2]} & \textbf{[A3]} \\
    w/o mask         & [A4] & [A5] & [A6] \\
    w/o prob (MSE)   & [A7] & [A8] & --- \\
    Shallow ($L=1$)  & [A9] & [A10] & [A11] \\
    \bottomrule
  \end{tabular}
\end{table}

\subsection{Diurnal Calibration}

Figure~\ref{fig:diurnal} shows the average predicted $\mu_h$ and $\pm
1\sigma_h$ intervals for a representative summer and winter day, averaged
across all 817 sites. REnFormer correctly assigns near-zero $\sigma$ at
peak irradiance hours and widens uncertainty at dawn/dusk transition
periods, where cloud-cover variability is highest. Per-site LSTM produces
calibrated means but cannot provide uncertainty estimates.

% -----------------------------------------------------------------------
\section{Discussion}
% -----------------------------------------------------------------------

\textbf{BESS co-location.}
The SEN Chile dataset includes 60 battery energy storage systems (BESS)
co-located with solar plants, with dispatch ranging from $-220$ to $+220$\,MW.
In this work we treat BESS as a grid-level phenomenon and do not include BESS
signals as input features. This is a pragmatic choice: BESS dispatch is
largely operator-controlled and therefore non-predictable from generation
history alone. Future work should investigate joint solar-BESS forecasting
as a multi-task problem.

\textbf{Distribution shift.}
Chile's solar fleet grew by approximately [X]\% between 2021 and 2026 as
measured by P99 peak generation. Our chronological split means the test
set contains newer, larger plants whose generation patterns were not
seen at training time. This represents realistic deployment conditions
and explains the moderate performance gap between our model and an
oracle upper bound.

\textbf{Limitations.}
REnFormer is trained without meteorological covariates (no NWP input),
relying solely on historical generation. Adding irradiance forecasts and
temperature would likely improve performance, particularly at the 6--24h
horizon. The current architecture also uses last-token pooling, which
may under-utilise mid-sequence context for long lookbacks.

% -----------------------------------------------------------------------
\section{Conclusion}
% -----------------------------------------------------------------------

We presented REnFormer, a global Transformer for probabilistic multi-site
solar power forecasting trained on 817 Chilean grid plants. The key
contributions are a parameter-efficient global architecture, a masked
Gaussian NLL loss that focuses learning on non-trivial generation steps,
and a large-scale evaluation on the most extensive real-grid solar dataset
in the global forecasting literature. REnFormer consistently outperforms
per-site LSTM and MLP, global TFT, and a MSE-trained ablation, while
providing calibrated uncertainty estimates throughout the diurnal cycle.
Code and pre-trained weights will be released upon acceptance.

% -----------------------------------------------------------------------
% References
% -----------------------------------------------------------------------
\bibliographystyle{IEEEtran}
\bibliography{renformer}

\end{document}
